\subsection{Fixed points: an additional necessary condition and a literature correction}

\begin{proposition}[Proved: preservation of monomorphisms of set diagrams]\label{D:set-monos}
If $A$ is totally $\Winfty$-aspherical, then
$\colim_A:\Set^A\to\Set$ preserves monomorphisms. Moreover every span
$a\xrightarrow{f}b$, $a\xrightarrow{g}c$ in $A$ has a commuting completion:
there are $u:b\to d$, $v:c\to d$ with $uf=vg$.
\end{proposition}
\begin{proof}
For a monomorphism $K\hookrightarrow H$, identify $K$ with a subfunctor.
If two objects of $\operatorname{El}_A(K)$ have the same image in
$\colim_AH$, they lie in one connected component $E$ of
$\operatorname{El}_A(H)$. Theorem~\ref{C:classification} makes $E$ totally
aspherical, and Lemma~\ref{C:directed-affine} gives a common successor of
these two objects in $E$. Since $\operatorname{El}_A(K)$ is a cosieve, this
successor and the two arrows belong to $\operatorname{El}_A(K)$.
Consequently the two objects also have the same image in $\colim_AK$.
This proves injectivity.

Apply the same common-successor assertion to
$\operatorname{El}_A(A(a,-))=a\backslash A$, which is connected because it
has an initial object. Its two objects $f$ and $g$ have a common successor
precisely when the displayed commuting completion exists.
\end{proof}

\begin{proposition}[Proved: fixed points in connected representations]\label{D:fixed-points}
Let $A$ be totally $\Winfty$-aspherical, and let $H:A\to\Set$ have a
nonempty connected category of elements. Every natural endomorphism
$\tau:H\to H$ fixes an element: there are $a\in A$ and $h\in H(a)$ with
$\tau_a(h)=h$.
\end{proposition}
\begin{proof}
We use Corollary~\ref{B:homotopical-colimits}: ordinary $A$-colimits of
simplicial sets preserve objectwise weak equivalences.

First suppose that no element of any $H(a)$ is periodic under $\tau_a$.
Let $X(a)$ be the nerve of the category of the $\mathbb N$-action on
$H(a)$ generated by $\tau_a$. An arrow $x\to y$ is an integer $m\geq0$
with $\tau_a^m(x)=y$; composition adds integers. In the absence of periodic
elements there is at most one arrow with given endpoints. Each connected
component is directed: elements joined by a finite zigzag of forward
iterates have a common forward iterate. Thus each component is filtered
and weakly contractible. The quotient map
\[
 X(a)\longrightarrow Q(a):=H(a)/\langle x\sim\tau_a(x)\rangle
\]
is an objectwise weak equivalence to a discrete simplicial set.

In simplicial degree $j$ one has $X(a)_j=H(a)\times\mathbb N^j$.
Product with the fixed set $\mathbb N^j$ preserves set colimits.
As $\colim_AH$ is a singleton, the face and degeneracy maps consequently
identify $\colim_AX$ with the nerve $B\mathbb N$ of the one-object monoid
$\mathbb N$. On the other hand, $\colim_AQ$ is a singleton. Homotopicality
of colimit would make $B\mathbb N$ weakly contractible, whereas its
edge-path presentation has fundamental group the group completion
$\mathbb Z$. This contradiction proves existence of a periodic element.

Let $r\geq1$ be the least positive period occurring anywhere, and put
\[
 P(a)=\{h\in H(a):\tau_a^r(h)=h\}.
\]
Every element of $P(a)$ has exact period $r$. Naturality makes $P$ a
subfunctor, and its category of elements is connected. Indeed, any two
of its objects have a common successor in the totally aspherical
category $\operatorname{El}_A(H)$; the successor belongs to $P$.
The restriction of $\tau$ therefore gives a free action of the cyclic
group $C_r$ on each $P(a)$.

The nerve $Y(a)$ of the action groupoid $P(a)//C_r$ maps by an objectwise
weak equivalence to the orbit set $P(a)/C_r$: each orbit groupoid has
trivial stabilizers and is equivalent to a point. Its degree-$j$ set is
$P(a)\times C_r^j$. Since $\colim_AP$ is a singleton, the colimits of
these two diagrams are $BC_r$ and a point. Homotopicality now makes
$BC_r$ weakly contractible, so its fundamental group $C_r$ is trivial.
Thus $r=1$, providing a fixed element.
\end{proof}

\begin{remark}[Known: the fixed-point condition does not imply exactness]\label{D:fpp-not-exact}
The fixed-point condition of Proposition~\ref{D:fixed-points}, for all
connected representations, is strictly weaker than exactness of
abelian colimits. A small connected counterexample is constructed in
\citet[Part~C, Construction~C.1 and Theorem~C.2, pp.~384--389]{AdamekKoubekReiterman1981Embeddings}.
Their Section~A.4, p.~370, explicitly corrects the earlier identification
of these conditions. In particular, the statement in
\citet[Section~II.3, p.~100]{AdamekReiterman1976SetColim} cannot be used as
an exactness criterion. The 1981 counterexample has the fixed-point
property but fails simultaneous affine equalization of three specified
parallel arrows.

This source correction does not resolve Remark~\ref{D:exact-converse-open}:
we have not proved that the published counterexample satisfies
$\mathcal P_\infty$. The distinction between pairwise and simultaneous
affine equalization remains essential.
\end{remark}
